% Supplementary results for Section~\ref{sec:confound}: the normalisation
% control built on the summed-detail variant. Input under the appendix section
% that carries \label{app:supp}, so this file starts at subsection level.

\subsection{A normalisation control on the summed-detail shape}
\label{app:supp-confound}

A further control is built on the summed-detail variant of
Eq.~\eqref{eq:summed}, whose per-band weights are
$(w_{\mathrm{LL}}, w_{\det}, w_{\det}, w_{\det})$, and fixes
$\sum_b \omega_b^2$ at \levelSqFixedSum, that variant's mean over the training
timesteps (Table~\ref{tab:confound-summed}). Because it differs from the
controls of Table~\ref{tab:confound} in band shape as well as in
normalisation, it is compared only with its parent, the unnormalised
summed-detail arm. Against that parent, the control changes FID by
$\dSqFixedSumVsEqSeven$ $[\dSqFixedSumVsEqSevenLo, \dSqFixedSumVsEqSevenHi]$
(raw $p = \dSqFixedSumVsEqSevenP$, $p_{\mathrm{Holm}} = \dSqFixedSumVsEqSevenHolm$
over the \nFamilyConfound{} confound contrasts), although it remains
worse than MSE ($+\dSqFixedSum$~FID $[\dSqFixedSumLo, \dSqFixedSumHi]$,
$p_{\mathrm{Holm}} = \dSqFixedSumHolm$). For this shape, whose total falls
from $3(1 - \alpha)$ to $\alpha$, holding $\sum_b \omega_b^2$ constant
therefore improves FID relative to the parent, whereas for the released,
detail-averaged shape the same normalisation raises FID relative to AWWL
(Section~\ref{sec:confound-results}). The effect of normalising the timestep
weight depends on the band shape it accompanies, and the controls support no
general statement that timestep reweighting helps or harms. As for the main
controls, the effective timestep weight $\Lambda(t)$ is not held fixed
(Proposition~\ref{prop:confound}(b)).

% Generated by scripts/make_tables.py --- do not edit by hand.
% Regenerate after every sweep; the paper reads these directly.

\begin{table}[tbp]
\centering
\footnotesize
\setlength{\tabcolsep}{3pt}
\caption{Normalisation control on the summed-detail variant of AWWL (each detail band carries $w_{\det}$), CIFAR-10, epoch 199, DDIM-100. The band weights are $\omega_b = (w_{\mathrm{LL}}, w_{\det}, w_{\det}, w_{\det})$; the control holds $\sum_b \omega_b^2$ at the variant's own mean over the training timesteps. $\Delta$ is the paired-by-seed FID difference with its 95\% $t$-interval over the $n$ seeds; $p$ is the two-sided paired $t$-test, raw and Holm-corrected. Against MSE the Holm family is the 21-arm family of Table~\ref{tab:compare-fid}; against the reference it is the four confound contrasts of Tables~\ref{tab:confound} and~\ref{tab:confound-summed}. Levels read from the run configurations: AWWL, detail bands summed, $\sum_b \omega_b^2$ fixed: $\sum_b \omega_b^2 = 0.2454$. Bold: $p_{\mathrm{Holm}} < 0.05$ within the stated family.}
\label{tab:confound-summed}
\begin{tabular}{>{\raggedright\arraybackslash\hangindent=1em\hangafter=1}p{3.3cm}rrrrrrr}
\toprule
 & & \multicolumn{3}{c}{Against MSE} & \multicolumn{3}{c}{Against its parent (first row)} \\
\cmidrule(lr){3-5}\cmidrule(lr){6-8}
Arm & $n$ & $\Delta$ [95\% CI] & $p$ & $p_{\mathrm{Holm}}$ & $\Delta$ [95\% CI] & $p$ & $p_{\mathrm{Holm}}$ \\
\midrule
AWWL, detail bands summed & 5 & $+2.63$\;$[2.22, 3.04]$ & $<10^{-4}$ & $\mathbf{0.0008}$ & \emph{reference} &  &  \\
AWWL, detail bands summed, $\sum_b \omega_b^2$ fixed & 5 & $+1.75$\;$[1.19, 2.32]$ & $0.0010$ & $\mathbf{0.0090}$ & $-0.87$\;$[-1.16, -0.58]$ & $0.0011$ & $\mathbf{0.0045}$ \\
\bottomrule
\end{tabular}
\end{table}

